The technological advancement of Unmanned Aerial Vehicles (UAVs) has significantly expanded their applicability, enabling them to overcome the mobility constraints of ground robots in tasks such as inspection, agriculture, and object transportation to inaccessible locations \cite{carvajal2024}. With this evolution, UAVs have ceased to function solely as observation platforms and have become capable of performing autonomous tasks involving physical interaction with the environment, such as pick-and-place missions (object capture and placement) \cite{abuzayed2020, cocuzza2021}. The autonomous execution of these tasks requires the establishment of control algorithms that enable the robot to operate safely \cite{carvajal2024}. Furthermore, in environments where global navigation systems (GPS) do not provide the precision required for aligning a robotic gripper, Visual Servoing emerges as the most viable solution, using feedback from onboard cameras to guide the aircraft \cite{shang2016}.

The autonomous capture of a non-cooperative target, such as a cage or a solid black block without visual markers, imposes significant challenges on classical perception methods and aircraft stability \cite{zhao2019}. The motivation for investigating this problem lies in the vulnerability of traditional vision systems when operating in outdoor and unstructured environments \cite{zhao2019}. In terms of perception, algorithms based exclusively on color space models, such as HSV, are highly sensitive to illumination variations. Likewise, classical geometric feature extraction techniques (such as SIFT, SURF, and HOG), combined with sliding window scanning strategies, require redundant computational effort and present difficulties in modeling the diversity of image backgrounds \cite{zhao2019}. As a consequence, these classical approaches frequently suffer from the “semantic gap,” demonstrating an inability to abstract scene context, which causes the system to confuse the absence of light from a dark object with shadows or noise from the environment itself \cite{zhao2019}.

Regarding control, Visual Servoing suffers from latency, since image processing is computationally intensive and operates at frequencies significantly lower than those of inertial flight sensors, resulting in a large sampling period and consequently generating intervals without visual feedback \cite{shang2016}. Furthermore, the capture strategy adopted in this project proposes that the UAV physically positions itself over the target, accommodating the cage between its landing gear so that a robotic gripper can be aligned with a specific capture handle \cite{fagundes2025}. This mechanical approach requires the visual system to provide not only the global positioning of the structure but also references for fine orientation alignment (yaw) during descent \cite{fagundes2025}.

Given the need to overcome these technological limitations, the following research question emerges: How can deep learning-based detection algorithms be integrated with dynamic control strategies to ensure the autonomous capture of a non-cooperative target by a UAV while mitigating illumination interference, image latency, and ensuring the physical alignment of the robotic gripper?

To address detection-related challenges, the literature has increasingly focused on Deep Learning models \cite{amjoud2023}. This is because the YOLO (You Only Look Once) family of Convolutional Neural Networks formulates detection as a single regression problem capable of predicting bounding boxes and probabilities in a single forward pass, ensuring high frame rates and accuracy \cite{afifah2025, amjoud2023, adegun2024}. The YOLOv8 architecture, in particular, introduced improvements such as the anchor-free design, which eliminates the need for predefined anchor boxes and demonstrates the ability to abstract deep semantic features to overcome illumination variations and occlusions \cite{afifah2025}.

However, although neural networks provide high accuracy in extracting semantic features, visual perception in UAVs is frequently affected by dynamic noise. Airborne object detection suffers from momentary failures and abrupt changes in the bounding box due to vibrations and occlusions caused by the propellers \cite{barisic2019}. To mitigate these failures and the noisy oscillations of neural networks applied to drones, the use of a Kalman Filter has been identified as a valid solution \cite{barisic2019}. This estimation technique ensures a continuous and reliable approximation not only of the target position but also of its kinematic velocity, compensating for deficiencies generated during periods of image tracking loss \cite{barisic2019}.

In addition to filtering the visual signal, the spatial reference frame adopted for navigation represents another challenge. This stems from the limitations of implementing the UAV control loop purely within the two-dimensional image plane (2D). When navigation is based solely on pixel coordinates, the system becomes susceptible to depth ambiguity and sensitive to geometric uncertainties, since altitude variations alter the scale of the object in the camera view. In this context, there are significant advantages to applying the analytical Pinhole camera model \cite{thuilot2002, fagundes2025}. The use of geometric transformation matrices enables the conversion of pixel coordinates directly into a three-dimensional (3D) Cartesian metric space, reducing sensitivity to uncertainties during the approach phase and providing substantially more stable flight dynamics \cite{fagundes2025}.

Based on this scenario, the main hypothesis of this work is that integrating an object detection algorithm for the semantic isolation of a non-cooperative target, operating in series with a Kalman Filter and a 2D–3D spatial transformation based on the Pinhole model, can mitigate visual noise and enable the stable flight control required for the execution of an autonomous object capture system using UAVs and Visual Servoing. Furthermore, it is hypothesized that, by supplying the control loop of a cascaded PD (Proportional-Derivative) Controller with the filtered velocity extracted from the estimator, the system will be capable of compensating for camera sampling delays and suppressing the overshoot caused by the traditional Integral term \cite{shang2016, jiang2021, fagundes2025}. Additionally, with the recent architectural evolution of detection networks, the need emerges to investigate the impact of new topologies on the optimization of this process. Therefore, this work proposes a comparative evaluation of different deep learning-based detection architectures in order to identify which model offers the best trade-off between accuracy and computational cost for the constrained task of onboard cage detection \cite{amjoud2023}.

Given the above, the general objective of this monograph is to develop and validate an autonomous object capture system using UAVs and Visual Servoing, aiming at the search, approach, and capture of a non-cooperative object in a simulated environment. To achieve the general objective, the following specific objectives were defined:

\begin{itemize}

\item Analyze deep learning algorithms for multi-stage target detection under illumination variations, aiming to align the robotic gripper with a Mean Average Precision (mAP) greater than 75\% and real-time processing above 20 frames per second (FPS);

\item Investigate the application of the Kalman Filter to mitigate perception noise and predict the target kinematics in the image plane, aiming to reduce the Root Mean Square Error (RMSE) of the estimated trajectory by at least 20\% when compared to the raw detection signal;

\item Evaluate the feasibility of mathematically converting the two-dimensional (2D) camera error into spatial coordinates (3D), establishing as a success criterion an absolute distance estimation error lower than 0.40 meters in the world Cartesian reference frame;

\item Propose and analyze the behavior of a flight control strategy with dynamic compensation (PD Controller) to overcome image latency inherent to Visual Servoing systems, ensuring a damped approach with trajectory overshoot below 15\%;

\item Validate the effectiveness of the proposed architecture through realistic physical simulation under the Robot Operating System (ROS), defining as a capture success criterion an approach success rate above 70\% and a final Euclidean distance lower than 0.30 meters between the UAV and the center of the cage.

\end{itemize}

The methodology of this research is characterized as a modeling and validation study conducted entirely in a simulation environment. For its execution, the software architecture is managed by the ROS (Robot Operating System) ecosystem, responsible for orchestrating message exchange within the Visual Servoing strategy, connecting computer vision algorithms (object detection and state estimation) to the flight control loops \cite{fagundes2025, sciortino2018}. The validation of UAV dynamics and sensor emulation is performed in the Gazebo realistic physics simulator, integrated with the flight controller through the Software-in-the-Loop (SITL) architecture. Communication between the ROS environment and the control firmware, ArduPilot or PX4, occurs through the MAVROS package, using the MAVLink protocol for the exchange of commands, states, and telemetry \cite{sciortino2018}.

Finally, the structural organization of this work is divided into five chapters. Chapter 1 comprises this Introduction, which presents the context, motivation, research question, and objectives of the study. Chapter 2 presents the Literature Review, segmented into two main parts: the Theoretical Foundation, which explores the essential concepts of Unmanned Aerial Vehicles (UAVs), Computer Vision Applied to UAVs, Visual Servoing, and Aerial Capture and Manipulation Systems; and the Related Works section, which discusses the state of the art in the literature. Chapter 3 details the Methodology, encompassing image collection, the implementation of deep learning algorithms for detection, and the structuring of the control logic. Chapter 4 presents the Results and Discussion, highlighting the performance comparison among the evaluated AI architectures, as well as the accuracy and behavior of the UAV in the simulator. Chapter 5 concludes with the study’s Conclusions and provides recommendations for future work.
